% =============================================================================
% 5.5 Flujo del Document-Driven SDLC (DDS)
% Archivo: 5.5_flujo_dds.tex
% =============================================================================

\section{Flujo del Document-Driven SDLC (DDS)}
\label{sec:flujo-dds}

El \textit{Document-Driven SDLC} (DDS) constituye la metodolog\'ia
propietaria adoptada para sistematizar el proceso completo de dise\~no,
desarrollo y documentaci\'on de la plataforma. A diferencia de las
metodolog\'ias \'agiles convencionales, el DDS integra expl\'icitamente las
fases de documentaci\'on t\'ecnica y validaci\'on acad\'emica dentro del flujo
de trabajo de ingenier\'ia de software, asegurando que cada artefacto producido
sea trazable, reproducible y verificable. El flujo se estructura en ocho fases
secuenciales, cada una con objetivos, entregables y criterios de aceptaci\'on
claramente definidos. El contraste sistem\'atico de esta metodolog\'ia frente a
los enfoques de cascada, \'agil/Scrum, RUP, \textit{dual-track} y de
\textit{design docs}/RFC ---junto con su fundamento acad\'emico e industrial---
se desarrolla en el marco te\'orico (\autoref{sec:flujo_dds_ia}) y se sintetiza en
la \autoref{tab:dds_comparativa}.

\subsection{Fase 0: Developer Experience (DX)}
\label{subsec:fase0-dx}

La Fase 0 establece las bases del entorno de desarrollo y la experiencia del
desarrollador (\textit{Developer Experience}, DX). Esta fase previa abarca la
configuraci\'on del entorno de desarrollo integrado (IDE), la selecci\'on y
estandarizaci\'on del \textit{stack} tecnol\'ogico, la definici\'on de
convenciones de c\'odigo (\textit{linting}, formateo, nomenclatura), y la
implementaci\'on de la infraestructura de integraci\'on y despliegue continuo
(CI/CD). El objetivo primordial es garantizar que todo miembro del equipo de
desarrollo pueda incorporarse al proyecto con fricci\'on m\'inima y productividad
inmediata.

\subsection{Fase 1: Descubrimiento y An\'alisis}
\label{subsec:fase1-descubrimiento}

La fase de descubrimiento y an\'alisis comprende el an\'alisis inverso del
sistema preexistente (\textit{As-Is}), la identificaci\'on de \textit{stakeholders}
y la derivaci\'on de los requisitos a partir del c\'odigo fuente y de los procesos
organizacionales reales. Los artefactos producidos en esta fase incluyen el
inventario del c\'odigo fuente, la documentaci\'on del estado actual, el mapa de
actores, los diagramas de contexto (C4) y el mapa de brechas identificadas.
Complementariamente, se ejecuta un ejercicio de modelado de amenazas
(\textit{Red Teaming} con STRIDE) previo a la escritura de c\'odigo nuevo, con lo
que la fase captura la realidad operativa del sistema antes de intervenirlo.

\subsection{Fase 2: Innovaci\'on y Validaci\'on}
\label{subsec:fase2-innovacion}

La fase de innovaci\'on y validaci\'on documenta los patrones de innovaci\'on
t\'ecnica del sistema y contrasta su viabilidad con las restricciones reales del
\textit{stack}. Entre estos patrones destacan el acceso a datos sin ORM
(\textit{Zero-ORM}) con \textit{Row Level Security} din\'amico, la arquitectura
h\'ibrida \textit{serverless} m\'as servidor \textit{stateful} de seguridad, y el
esquema \textit{Multi-Page Application} nativo sobre Vite. La validaci\'on se
formaliza mediante una matriz de integraci\'on que traza cada requisito hacia su
implementaci\'on t\'ecnica y registra su estado y sus riesgos, permitiendo
descartar aproximaciones inviables antes de escalar la construcci\'on. Los
hallazgos de esta fase alimentan directamente la especificaci\'on de requisitos
funcionales y no funcionales del sistema.

\subsection{Fase 3: Dise\~no y Definici\'on}
\label{subsec:fase3-diseno}

La fase de dise\~no y definici\'on transforma los requisitos validados en
especificaciones t\'ecnicas detalladas. Se elaboran los diagramas de casos de
uso, los diagramas de secuencia, los modelos de datos y los contratos de la
API REST. Adicionalmente, se definen las pol\'iticas de seguridad, los
esquemas de autorizaci\'on basados en roles (RBAC) y las reglas de negocio
formalizadas. Esta fase produce el \textit{blueprint} arquitect\'onico que
guiar\'a las fases subsiguientes de implementaci\'on.

\subsection{Fase 4: UI/UX}
\label{subsec:fase4-uiux}

La fase de dise\~no de interfaz de usuario (UI) y experiencia de usuario (UX)
abarca la creaci\'on del sistema de dise\~no visual, la elaboraci\'on de
\textit{wireframes} de baja y alta fidelidad, y la construcci\'on de prototipos
interactivos navegables. Se aplican principios de dise\~no centrado en el
usuario, heur\'isticas de usabilidad de Nielsen y directrices de accesibilidad
WCAG 2.1. Los prototipos resultantes son evaluados mediante pruebas de
usabilidad con usuarios finales, cuyos hallazgos retroalimentan el dise\~no
antes de la implementaci\'on del \textit{frontend}.

\subsection{Fase 5: Arquitectura y Desarrollo}
\label{subsec:fase5-arquitectura}

La fase de arquitectura y desarrollo materializa las especificaciones t\'ecnicas
en c\'odigo fuente funcional. Se implementa la arquitectura multi-inquilino
(\textit{multi-tenant}), los microservicios de autenticaci\'on y autorizaci\'on,
el motor de evaluaci\'on de riesgos, y los m\'odulos de gesti\'on operativa.
El desarrollo sigue un enfoque de \textit{Test-Driven Development} (TDD),
donde los casos de prueba se escriben previamente a la implementaci\'on de la
funcionalidad, garantizando una cobertura de c\'odigo elevada desde las etapas
iniciales del desarrollo.

\subsection{Fase 6: QA y Testing}
\label{subsec:fase6-qa}

La fase de aseguramiento de calidad (QA) y pruebas comprende la ejecuci\'on
sistem\'atica de pruebas unitarias, pruebas de integraci\'on, pruebas de la
API REST y pruebas de seguridad. Se emplean herramientas automatizadas como
Jest para las pruebas unitarias del \textit{backend}, y se implementan
escenarios de prueba que cubren tanto los flujos principales como los flujos
alternativos y de error. Los resultados de cobertura son analizados
cuantitativamente para identificar \'areas de riesgo y priorizar esfuerzos
de prueba adicionales.

\subsection{Fase 7: Despliegue y Operaciones}
\label{subsec:fase7-despliegue}

La fase final de despliegue y operaciones abarca la configuraci\'on de la
infraestructura en la nube, la implementaci\'on de \textit{pipelines} de
despliegue continuo, la configuraci\'on de monitoreo y alertas, y la
elaboraci\'on de la documentaci\'on operativa. Se definen los procedimientos
de \textit{rollback}, las estrategias de escalamiento horizontal y las
pol\'iticas de respaldo y recuperaci\'on ante desastres. Esta fase asegura
que el sistema sea operable, observable y mantenible en un entorno de
producci\'on real.

\subsection{Diagrama del Flujo DDS}
\label{subsec:diagrama-flujo-dds}

La Figura~\ref{fig:flujo-dds} presenta el diagrama visual del flujo completo
del \textit{Document-Driven SDLC}, ilustrando la secuencia de fases,
sus interdependencias y los puntos de retroalimentaci\'on entre etapas.

\begin{figure}[htbp]
\centering
\includegraphics[width=0.95\textwidth]{Flujo de DDS}
\caption{Diagrama del flujo del \textit{Document-Driven SDLC} (DDS) con sus ocho fases}
\label{fig:flujo-dds}
\end{figure}

\subsection{Mapeo de Fases y Artefactos}
\label{subsec:mapeo-fases-artefactos}

La Tabla~\ref{tab:fases-artefactos} establece la correspondencia entre cada
fase del DDS y los artefactos tangibles producidos, proporcionando trazabilidad
completa del proceso de documentaci\'on.

\begin{small}
\begin{longtable}{|p{1cm}|p{3.2cm}|p{5cm}|p{4.5cm}|}
\caption{Mapeo de fases del DDS con artefactos y entregables}
\label{tab:fases-artefactos} \\
\hline
\textbf{Fase} & \textbf{Nombre} & \textbf{Artefactos principales} & \textbf{Criterio de aceptaci\'on} \\
\hline
\endfirsthead

\multicolumn{4}{c}%
{{\tablename\ \thetable{} -- Continuaci\'on de la p\'agina anterior}} \\
\hline
\textbf{Fase} & \textbf{Nombre} & \textbf{Artefactos principales} & \textbf{Criterio de aceptaci\'on} \\
\hline
\endhead

\hline
\multicolumn{4}{r}{{Contin\'ua en la siguiente p\'agina\ldots}} \\
\endfoot

\hline
\endlastfoot

0 & Developer Experience & Configuraci\'on de IDE, convenciones de c\'odigo, \textit{pipelines} CI/CD & Entorno reproducible y documentado \\
\hline
1 & Descubrimiento y An\'alisis & Inventario de c\'odigo, an\'alisis \textit{As-Is}, mapa de actores, diagramas C4, mapa de brechas, \textit{Red Teaming} (STRIDE) & Requisitos derivados y trazados al sistema real \\
\hline
2 & Innovaci\'on y Validaci\'on & Patrones de innovaci\'on (Zero-ORM+RLS, arquitectura h\'ibrida, MPA), matriz de integraci\'on de requisitos & Viabilidad t\'ecnica confirmada sobre el \textit{stack} \\
\hline
3 & Dise\~no y Definici\'on & Casos de uso, diagramas de secuencia, modelo de datos, contratos API & Especificaciones t\'ecnicas aprobadas \\
\hline
4 & UI/UX & Sistema de dise\~no, \textit{wireframes}, prototipos interactivos & Pruebas de usabilidad superadas \\
\hline
5 & Arquitectura y Desarrollo & C\'odigo fuente, pruebas TDD, documentaci\'on t\'ecnica & Funcionalidad implementada y verificada \\
\hline
6 & QA y Testing & Reportes de cobertura, matrices de trazabilidad, informes de defectos & Cobertura $\geq$ 80\%, cero defectos cr\'iticos \\
\hline
7 & Despliegue y Operaciones & Infraestructura configurada, monitoreo activo, documentaci\'on operativa & Sistema desplegado y observable \\

\end{longtable}
\end{small}

\subsection{Arquitectura de la Plataforma}
\label{subsec:arquitectura-plataforma}

La implementaci\'on del flujo DDS se sustenta en una arquitectura de plataforma
cuidadosamente dise\~nada para soportar las necesidades operativas de
m\'ultiples organizaciones no gubernamentales de manera simult\'anea. La
Figura~\ref{fig:arquitectura-plataforma} presenta el diagrama de arquitectura
general de la plataforma, donde se visualizan los componentes principales, sus
interacciones y los flujos de datos entre las capas del sistema.

\begin{figure}[htbp]
\centering
\includegraphics[width=0.95\textwidth]{Diagrama de arquitectura de la plataforma de ONG}
\caption{Diagrama de arquitectura de la plataforma de gesti\'on de ONG}
\label{fig:arquitectura-plataforma}
\end{figure}

La arquitectura adopta un patr\'on de capas claramente diferenciadas: la capa
de presentaci\'on (\textit{frontend}) implementada en React con TypeScript, la
capa de l\'ogica de negocio (\textit{backend}) construida sobre Node.js con
Express, y la capa de persistencia basada en PostgreSQL con aislamiento de
datos por esquema para garantizar la segregaci\'on multi-inquilino. Los
servicios transversales de seguridad, auditor\'ia e inteligencia artificial
operan como componentes desacoplados que pueden escalar de forma independiente.

\subsection{Stack Tecnol\'ogico y Estandarizaci\'on Arquitect\'onica}
\label{subsec:stack-tecnologico}

La selecci\'on del \textit{stack} tecnol\'ogico responde a criterios de madurez,
rendimiento, ecosistema de herramientas y alineaci\'on con las competencias del
equipo de desarrollo. La Figura~\ref{fig:stack-tecnologico} presenta el diagrama
estructural del \textit{stack} tecnol\'ogico adoptado y la estandarizaci\'on
arquitect\'onica aplicada a todos los m\'odulos del sistema.

\begin{figure}[htbp]
\centering
\includegraphics[width=0.95\textwidth]{Diagrama estructural del stack tecnologico y estandarizacion arquitectonica del sistema}
\caption{Diagrama estructural del \textit{stack} tecnol\'ogico y estandarizaci\'on arquitect\'onica del sistema}
\label{fig:stack-tecnologico}
\end{figure}

La estandarizaci\'on arquitect\'onica garantiza que todos los m\'odulos del
sistema sigan patrones de dise\~no consistentes, lo cual reduce la curva de
aprendizaje para nuevos desarrolladores, facilita las revisiones de c\'odigo
y minimiza la deuda t\'ecnica acumulada. Los est\'andares definidos incluyen
la estructura de directorios, las convenciones de nomenclatura, los patrones
de manejo de errores, las estrategias de \textit{logging} y los esquemas de
validaci\'on de datos de entrada.
